% ============================================================
% conclusion.tex — Kết luận
% ============================================================

\chapter*{KẾT LUẬN}
\addcontentsline{toc}{chapter}{KẾT LUẬN}
\markboth{KẾT LUẬN}{KẾT LUẬN}

% ---------------------------------------------------------------
\section*{1. Tóm tắt}
\addcontentsline{toc}{section}{1. Tóm tắt}
% ---------------------------------------------------------------

Đề cương nghiên cứu đã trình bày bài toán, phương pháp và kế hoạch thực hiện cho đề tài \textit{``\detai''}. Đề tài hướng đến xây dựng một prototype nghiên cứu end-to-end cho bài toán dự đoán lan truyền khói cấp phòng/khu vực trong tòa nhà cao tầng, tích hợp các thành phần:

\begin{itemize}
  \item Module BIM/IFC-to-Graph chuyển đổi tự động mô hình tòa nhà thành đồ thị.
  \item Mô phỏng batch bằng CFAST kết hợp đối sánh chọn lọc với FDS.
  \item Hệ thống IoT giả lập 3 lớp tạo dữ liệu cảm biến có nhiễu, trễ, mất dữ liệu và lỗi mô phỏng.
  \item Mô hình AI surrogate dạng Temporal GNN nhận cấu hình kịch bản đã biết và lịch sử IoT $[t-k,t]$ để dự báo \texttt{smoke\_label} tại $t+\Delta$ và \texttt{time\_to\_arrival}.
  \item Dashboard trực quan hóa và replay kết quả.
\end{itemize}

% ---------------------------------------------------------------
\section*{2. Tính mới và đóng góp dự kiến}
\addcontentsline{toc}{section}{2. Tính mới và đóng góp dự kiến}
% ---------------------------------------------------------------

Tính mới của đề tài nằm ở \textbf{pipeline tích hợp}, không phải ở việc phát minh riêng lẻ từng thành phần. Cụ thể:

\begin{enumerate}
  \item \textbf{Pipeline BIM → Simulation → AI:} Kết nối liền mạch từ mô hình BIM đến dự đoán AI, cho phép tự động hóa toàn bộ quy trình.
  \item \textbf{Building graph từ BIM/IFC:} Tận dụng thông tin kiến trúc có sẵn thay vì xây dựng mô hình thủ công.
  \item \textbf{Virtual IoT 3 lớp:} Tạo dữ liệu huấn luyện AI phản ánh tính bất hoàn hảo của cảm biến thực tế.
  \item \textbf{Boundary-Aware Temporal GNN:} Mô hình AI tích hợp cấu trúc đồ thị, lịch sử IoT, điều kiện biên và trạng thái hệ thống đã quan sát được đến thời điểm dự báo.
  \item \textbf{Đánh giá đa chiều:} So sánh BAT-GNN với ba nhóm baseline bắt buộc, thực hiện scenario/fire-node/fire-floor split, ablation và kiểm tra robustness.
\end{enumerate}

% ---------------------------------------------------------------
\section*{3. Giới hạn nghiên cứu}
\addcontentsline{toc}{section}{3. Giới hạn nghiên cứu}
% ---------------------------------------------------------------

Nghiên cứu này là một prototype thử nghiệm học thuật và tồn tại một số giới hạn cốt lõi: hoàn toàn sử dụng dữ liệu mô phỏng cháy; CFAST là nguồn sinh dữ liệu huấn luyện và FDS chỉ là tập đối sánh số chọn lọc; hệ thống cảm biến IoT được giả lập; mô hình biết trước cấu hình kịch bản; và phạm vi một tòa nhà chính không đủ để kết luận tổng quát hóa sang tòa nhà mới. Chi tiết được trình bày tại Chương~\ref{chap:assumptions}.

% ---------------------------------------------------------------
\section*{4. Hướng phát triển tương lai}
\addcontentsline{toc}{section}{4. Hướng phát triển tương lai}
% ---------------------------------------------------------------

Các nội dung dưới đây nằm ngoài phạm vi và tiêu chí nghiệm thu của đề tài hiện tại; chúng chỉ là gợi ý cho nghiên cứu tiếp theo sau khi hoàn thành P0.

\begin{itemize}
  \item Mở rộng số case đối sánh FDS lên 8--12 nếu đủ phần cứng.
  \item Tích hợp IoT phần cứng thật (cảm biến khói, nhiệt, CO).
  \item Mô hình hóa chi tiết hệ thống HVAC (duct, damper).
  \item Tích hợp mô hình thoát nạn (evacuation) để tính ASET/RSET.
  \item Mô phỏng CFD chi tiết cho atrium và không gian lớn.
  \item Phát triển hệ thống điều khiển PCCC thông minh dựa trên dự đoán AI.
  \item Kiểm chứng trên nhiều tòa nhà thực tế để đánh giá tính tổng quát.
\end{itemize}
